\section{Reference Implementation Notes}
\label{sec:impl}
% ================================================================================

This section sketches how a generic agent-runtime bridge can implement
the Bridge Contract without requiring changes to AGTP, GAR, or PACE.
The implementation is on the bridge side only.

\paragraph{State store.} The bridge maintains a state store keyed by
\texttt{Session-ID}, mapping each session to its current state, its
Session Identity Descriptor, its tool queue, and its event log. The store
is private to the bridge; AGTP, GAR, and PACE components do not read
from it directly.

\paragraph{Event emitter.} The bridge emits events in a JSON shape that
is compatible with GAR's SAR record (causal ordering, hash chaining).
In a GEC-present deployment, the events are forwarded to the GEC for
signing; the GEC remains the sole authoritative source. In a GEC-less
deployment (Definition~\ref{def:gecless}), the bridge SHOULD sign the
events itself with an operator key whose JWKS endpoint is published
through an out-of-band channel bound to the operator's deployment
identity (GAR \S 6.1, Invariants I-SIGN-EVENT and I-SIGN-JWKS). The bridge MUST NOT accept a
JWKS endpoint learned from session-supplied metadata.

\paragraph{Challenge store.} The bridge keeps, per session, a small
volatile set of live challenge nonces with their issue times and a
used flag (I9, I11). It is not durable: a restart voids every nonce and
clients ask again. The registry of client keys is scoped to a session and is rebuilt from the
log, which records each key's scheme tag and public key, and a restart
compares the stored registry with it.
Everything else needed to resume, the chain and the resume record, is in
the log (I12). The Python reference bridge and the Rust crate verify
resume requests signed with Ed25519 or ECDSA P-256 (clients and fixtures
sign), and their key registries
hold a scheme tag and public keys only, so the evidence exercises proof of possession against verification
material alone.

\paragraph{Tool-call queue.} While the human dimension is pending,
the bridge maintains a bounded queue of pending tool calls, each with
an idempotency key. Each enqueued call emits
\texttt{TOOL\_CALL\_QUEUED}. Once the client is attached and no
escalation is outstanding, the queue is drained in arrival order, each
call re-validated against the current scope and approval (I13).

\paragraph{Durable log and recovery.} The event log is written with a
commit barrier before any state change or acknowledgement (I12). On
start the bridge replays each log, verifies the chain, and recovers
sessions as detached (Section~\ref{sec:recovery}).

\paragraph{Findings from integration.} Mapping the contract onto a real
product (Termenbox ACP workspaces: the server is the bridge, the browser
the client, the agent a child process on stdio) found, as reported by the
author (the product is closed source and no reader can check this), problems that the
paper's own evidence could not, because the evidence came from the same
author reading the same text. A resume token minted at detach and
superseded by every later change while the client was away, could not
reach a client that vanished abruptly, and I9 then terminated the
session when it expired; the model had one client where a product owner
has several devices attached at once; the bridge sees an approval
request, not tool execution, so I2 as first worded promised more than a
bridge can see; an agent announces its own mode changes, which a bridge
can observe and revert but not prevent; and the independent defaults for
the human and call timers ended work that a mobile owner expects to
survive. Each, as reported, led to a change of the specification, not of the product
(challenge-response resume, a client set, answer binding, the
coverage limits of I2 and I3, derived timers). The same author wrote
both sides, which limits what this shows.

\paragraph{Reference implementation.} A reference implementation is
provided in the conformance test suite (Section~\ref{sec:conformance}).
The suite includes a \texttt{Bridge} class that the suite and its mutants
exercise against invariants $I_1, \ldots, I_{15}$; known deviations are
recorded in the limitations register (Table~\ref{tab:limitations}). It also has a
\texttt{SessionIdentity}, \texttt{Event}, and re-attach request data
model. A second implementation in Rust
(\texttt{supplementary/rust-bridge/}), written from Sections~4 and~5
rather than ported from the Python code (though it shares comments, error
texts and the key format with it), runs the same JSON vectors except
V52--V54, V88 and V89, which expect checkpoint events the crate does not
emit; it
has the same author, so it shows consistency of one reading of the
text. A runtime such as
Termenbox~\cite{termenbox}, a self-hosted web terminal for AI coding agents
(the author's project; see the note on evidence at the front of this
document), is one
example of a runtime that could implement the Bridge Contract; the
contract is deliberately implementation-agnostic and does not depend
on any particular runtime substrate.

% ================================================================================
